\uuid{mK7w}
\titre{Perceptron multicouche et détection d'un triangle du plan}
\niveau{L3}
\module{Analyse de données}
\chapitre{Réseaux de neurones}
\sousChapitre{Architecture et propagation avant}
\theme{Fonction d'activation Heaviside, séparation linéaire, région convexe}
\auteur{Maxime Nguyen}
\datecreate{2026-03-30}
\organisation{AMSCC}
\difficulte{2}

\contenu{
	\texte{
		On considère le triangle $T$ du plan dont les sommets sont
		$A = (1, 1)$, $B = (3, 1)$ et $C = (2, 3)$.
		On souhaite construire un réseau de neurones à propagation avant qui implémente
		la fonction indicatrice $f : \mathbb{R}^2 \to \{0, 1\}$ définie par
		\[
		f(x, y) = \begin{cases} 1 & \text{si } (x,y) \in T, \\ 0 & \text{sinon.} \end{cases}
		\]
		Chaque neurone dispose d'une fonction d'activation \textbf{Heaviside}
		$H : \mathbb{R} \to \{0,1\}$ définie par
		\[
		H(z) = \begin{cases} 1 & \text{si } z \geq 0, \\ 0 & \text{sinon,} \end{cases}
		\]
		et calcule donc, pour des poids $(w_1, w_2)$, un biais $b$ et des entrées $(x,y)$ :
		\[
		H(w_1 x + w_2 y + b).
		\]
	}
	
	\begin{enumerate}
		
		\item
		\question{
			Déterminer les équations des trois droites $(AB)$, $(BC)$ et $(CA)$.
			Pour chacune, préciser de quel côté se trouve l'intérieur du triangle,
			et écrire une inégalité de la forme $w_1 x + w_2 y + b \geq 0$
			qui caractérise le demi-plan contenant $T$.
		}
		\indication{
			Vérifier le signe de chaque expression en un point intérieur évident,
			par exemple le barycentre $G = (2, \frac{5}{3})$.
		}
		\reponse{
			\textbf{Droite $(AB)$} : $A=(1,1)$, $B=(3,1)$, donc $(AB)$ est la droite $y = 1$.
			$C = (2,3)$ vérifie $y - 1 = 2 > 0$, donc le demi-plan contenant $T$ est
			\[ y - 1 \geq 0. \]
			Neurone $n_1$ : $w_1 = 0$, $w_2 = 1$, $b = -1$.
			
			\medskip
			\textbf{Droite $(BC)$} : $B=(3,1)$, $C=(2,3)$.
			Vecteur directeur $\overrightarrow{BC} = (-1, 2)$, donc un vecteur normal est $(2, 1)$.
			Équation : $2(x-3) + 1(y-1) = 0$, soit $2x + y - 7 = 0$.
			$A=(1,1)$ donne $2(1)+1-7 = -4 < 0$, donc le demi-plan contenant $T$ est
			\[ -2x - y + 7 \geq 0. \]
			Neurone $n_2$ : $w_1 = -2$, $w_2 = -1$, $b = 7$.
			
			\medskip
			\textbf{Droite $(CA)$} : $C=(2,3)$, $A=(1,1)$.
			Vecteur directeur $\overrightarrow{CA} = (-1,-2)$, donc un vecteur normal est $(2,-1)$.
			Équation : $2(x-1) - 1(y-1) = 0$, soit $2x - y - 1 = 0$.
			$B=(3,1)$ donne $2(3)-1-1 = 4 > 0$, donc le demi-plan contenant $T$ est
			\[ 2x - y - 1 \geq 0. \]
			Neurone $n_3$ : $w_1 = 2$, $w_2 = -1$, $b = -1$.
		}
		
		\item
		\question{
			Montrer que $(x,y) \in T$ si et seulement si
			$n_1(x,y) + n_2(x,y) + n_3(x,y) = 3$.
			En déduire les poids, biais et activation d'un neurone de sortie $n_s$,
			prenant en entrées $n_1, n_2, n_3$, tel que $n_s(x,y) = f(x,y)$.
		}
		\reponse{
			Le triangle $T$ est l'intersection des trois demi-plans définis à la question
			précédente. Chaque $n_i$ vaut $1$ si le point est dans le $i$-ème demi-plan,
			$0$ sinon. Donc $n_1 + n_2 + n_3 = 3$ si et seulement si les trois contraintes
			sont satisfaites, c'est-à-dire si et seulement si $(x,y) \in T$.
			
			On pose :
			\[
			n_s = H(n_1 + n_2 + n_3 - 3),
			\]
			neurone de poids $(1,1,1)$ sur $(n_1,n_2,n_3)$ et de biais $-3$.
			On a $n_1+n_2+n_3-3 \geq 0 \iff n_1+n_2+n_3=3 \iff (x,y)\in T$,
			donc $n_s(x,y) = f(x,y)$.
		}
		
		\item
		\question{
			Décrire l'architecture complète du réseau : nombre de couches, nombre de neurones
			par couche. Récapituler dans un tableau tous les poids et biais.
			Vérifier ensuite que le réseau produit la bonne sortie pour les points
			$P_1 = (2, 2)$, $P_2 = (3, 3)$ et $P_3 = (1, 1)$.
		}
		\reponse{
			Le réseau comporte :
			\begin{itemize}
				\item une \textbf{couche d'entrée} : 2 neurones ($x$ et $y$) ;
				\item une \textbf{couche cachée} : 3 neurones $n_1, n_2, n_3$ (activation Heaviside) ;
				\item une \textbf{couche de sortie} : 1 neurone $n_s$ (activation Heaviside).
			\end{itemize}
			
			\begin{center}
				\renewcommand{\arraystretch}{1.3}
				\begin{tabular}{lcccc}
					\hline
					Neurone & Entrées & $w_x$ ou $w_{n_1}$ & $w_y$ ou $w_{n_2}$ & biais $b$ \\
					\hline
					$n_1$ & $x,y$ & $0$  & $1$  & $-1$ \\
					$n_2$ & $x,y$ & $-2$ & $-1$ & $7$  \\
					$n_3$ & $x,y$ & $2$  & $-1$ & $-1$ \\
					$n_s$ & $n_1,n_2,n_3$ & $1$ & $1$ & $-3$ (+ $w_{n_3}=1$) \\
					\hline
				\end{tabular}
			\end{center}
			
			\textbf{Vérification pour $P_1 = (2,2)$} (intérieur du triangle) :
			\begin{align*}
				n_1 &= H(2-1) = H(1) = 1, \\
				n_2 &= H(-4-2+7) = H(1) = 1, \\
				n_3 &= H(4-2-1) = H(1) = 1, \\
				n_s &= H(1+1+1-3) = H(0) = 1. \quad \checkmark
			\end{align*}
			
			\textbf{Vérification pour $P_2 = (3,3)$} (extérieur, au-delà de $(BC)$) :
			\begin{align*}
				n_1 &= H(3-1) = 1, \\
				n_2 &= H(-6-3+7) = H(-2) = 0, \\
				n_s &= H(1+0+\cdots - 3) \leq H(-1) = 0. \quad \checkmark
			\end{align*}
			
			\textbf{Vérification pour $P_3 = (1,1)$} (sommet $A$, sur la frontière) :
			\begin{align*}
				n_1 &= H(1-1) = H(0) = 1, \\
				n_2 &= H(-2-1+7) = H(4) = 1, \\
				n_3 &= H(2-1-1) = H(0) = 1, \\
				n_s &= H(3-3) = H(0) = 1. \quad \checkmark
			\end{align*}
		}
		
		\item
		\question{
			On suppose maintenant que l'on souhaite construire un réseau détectant l'union
			de deux triangles disjoints $T_1$ et $T_2$ du plan.
			Proposer une architecture (sans calculer les poids) permettant d'implémenter
			la fonction indicatrice de $T_1 \cup T_2$, et justifier pourquoi
			une seule couche cachée ne suffit pas en général pour une telle région.
		}
		\reponse{
			\textbf{Architecture proposée :}
			On ajoute une deuxième couche cachée.
			\begin{itemize}
				\item \textit{Première couche cachée} : 6 neurones, trois pour chaque triangle
				($n_1^{(1)}, n_2^{(1)}, n_3^{(1)}$ pour $T_1$ et
				$n_1^{(2)}, n_2^{(2)}, n_3^{(2)}$ pour $T_2$).
				\item \textit{Deuxième couche cachée} : 2 neurones $s_1$ et $s_2$,
				chacun implémentant la détection d'un triangle comme à la question
				précédente ($s_i = H(n_1^{(i)}+n_2^{(i)}+n_3^{(i)}-3)$).
				\item \textit{Couche de sortie} : 1 neurone $n_s = H(s_1 + s_2 - 1)$,
				qui vaut $1$ dès que l'un des deux triangles détecte le point.
			\end{itemize}
			
			\textbf{Pourquoi une seule couche cachée ne suffit pas :}
			Une seule couche cachée de neurones Heaviside ne peut implémenter que des
			intersections de demi-plans (régions convexes). Or $T_1 \cup T_2$ est en général
			une région \emph{non convexe} dès que les deux triangles sont disjoints :
			il existe des points sur le segment reliant un point de $T_1$ à un point de $T_2$
			qui n'appartiennent ni à $T_1$ ni à $T_2$, et qu'aucune intersection de
			demi-plans ne peut exclure tout en incluant les deux triangles.
			Une deuxième couche est donc nécessaire pour combiner les deux détecteurs convexes
			par une opération de type \og ou logique \fg.
		}
		
	\end{enumerate}
}